
\noindent\emph{Catatan.} Penting untuk diingat bahwa kita hanya bekerja dengan aljabar
\emph{kompleks}. Hasil ini \emph{salah} untuk aljabar Banach real. Contoh
tandingan sederhana adalah matriks $\begin{bmatrix}0 & 1 \\ -1 & 0 \end{bmatrix}$ yang dipandang sebagai
elemen aljabar Banach (real) $M_2$ yang terdiri atas semua matriks $2 \times 2$ dengan entri
bilangan real.

 \vskip 1 pt

\begin{thm}[Gelfand-Mazur]\label{C073141} Jika $A$ aljabar Banach beridentitas yang setiap
 \index{teorema Gelfand-Mazur}%
elemen taknolnya invertibel, maka $A$ isomorfik secara isometrik dengan~$\C$.
\end{thm}

\begin{proof}[\emph{Petunjuk untuk bukti}] Misalkan $B = \{ \lambda\vc 1\colon \lambda \in \C\}$. Gunakan
proposisi sebelumnya~\ref{C073134} untuk menunjukkan bahwa $B = A$.    \ns
\end{proof}

Pernyataan berikut merupakan konsekuensi langsung dari \emph{teorema Gelfand-Mazur}
\ref{C073141} dan proposisi~\ref{C037821}.

\begin{cor}\label{C073147} Jika $I$ ideal maksimal dalam aljabar Banach komutatif beridentitas $A$,
maka $A/I$ isomorfik secara isometrik dengan~$\C$.
\end{cor}

\begin{prop}\label{0006703fa} Misalkan $a$ suatu elemen dari aljabar beridentitas. Maka
   \[ \sigma(a^n) = [\sigma(a)]^n \]
untuk setiap $n \in \N$. (Notasi $[\sigma(a)]^n$ berarti $\{\lambda^n\colon \lambda \in
\sigma(a)\}$.)
\end{prop}

\begin{defn}\label{C073154} Misalkan $a$ suatu elemen dari aljabar beridentitas.
 \index{spektral!radius}%
 \index{radius!spektral}%
\df{radius spektral} dari~$a$, yang dinotasikan dengan $\rho(a)$, didefinisikan sebagai
$\sup\{\abs\lambda\colon \lambda \in \sigma(a)\}$.
\end{defn}

\begin{prop}\label{0006703} Jika $a$ suatu elemen dari aljabar Banach beridentitas, maka $\rho(a) \le
\norm a$ dan $\rho(a^n) = \bigl(\rho(a)\bigr)^n$.
\end{prop}

\begin{prop} Misalkan $\phi\colon A \sto B$ homomorfisme yang mempertahankan identitas di antara aljabar-aljabar beridentitas. Jika $a$ elemen invertibel dari $A$,
maka $\phi(a)$ invertibel dalam $B$ dan inversnya adalah~$\phi(a^{-1})$.
\end{prop}

\begin{cor} Jika $\phi\colon A \sto B$ homomorfisme yang mempertahankan identitas di antara aljabar-aljabar beridentitas, maka
   \[ \sigma(\phi(a)) \subseteq \sigma(a) \]
untuk setiap $a \in A$.
\end{cor}

\begin{cor} Jika $\phi\colon A \sto B$ homomorfisme yang mempertahankan identitas di antara aljabar-aljabar beridentitas, maka
   \[ \rho(\phi(a)) \le \rho(a) \]
untuk setiap $a \in A$.
\end{cor}

\begin{thm}[Rumus radius spektral]\label{C073157} Jika $a$ suatu elemen dari aljabar Banach
 \index{spektral!radius!rumus untuk}%
beridentitas, maka
   \[ \rho(a) = \inf\bigl\{\norm{a^n}^{1/n}\colon n \in \N\bigr\}
                                = \lim_{n \sto \infty} \norm{a^n}^{1/n}\,. \]
\end{thm}

\begin{exer} Dengan menggunakan rumus untuk operator Volterra $V$ dalam contoh~\ref{000319} sebagai
 \index{operator Volterra!dan persamaan integral}%
 \index{operator!Volterra!dan persamaan integral}%
acuan, pada ruang Banach~$\fml C([0,1])$ definisikan operator integral melalui rumus $Vf(x)=\int_0^x f(t)\,dt$.
 \begin{enumerate}
  \item[(a)] Hitung spektrum operator $V$. (\emph{Petunjuk.} Tunjukkan bahwa $\norm{V^n}
\le (n!)^{-1}$ untuk setiap $n \in \N$ dan gunakan \emph{rumus radius spektral}.)
  \item[(b)] Apa implikasi (a) terhadap kemungkinan memperoleh fungsi kontinu $f$ yang
memenuhi persamaan integral
   \[ f(x) - \mu\int_0^xf(t)\,dt = h(x), \]
dengan $\mu$ suatu skalar dan $h$ suatu fungsi kontinu yang diberikan pada~$[0,1]$?
  \item[(c)] Gunakan gagasan dalam (b) dan uraian deret Neumann untuk $\vc 1 - \mu V$ (lihat
proposisi~\ref{prop_Neumann_series}) untuk menghitung secara eksplisit (dan dalam bentuk fungsi elementer)
suatu solusi untuk persamaan integral
 \[f(x) - {\textstyle \frac12}\int_0^xf(t)\,dt = e^x.\]
 \end{enumerate}
\end{exer}

\begin{prop} Misalkan $A$ aljabar Banach beridentitas dan $B$ subaljabar tertutup yang memuat~$\vc 1_A$.
Maka
 \begin{enumerate}[label=\rm{(\alph*)}]
   \item $\inv B$ terbuka sekaligus tertutup dalam $B \cap \inv A$;
   \item $\sigma_A(b) \subseteq \sigma_B(b)$ untuk setiap $b \in B$; dan
   \item jika $b \in B$ dan $\sigma_A(b)$ tidak memiliki lubang (yakni, jika komplemennya dalam $\C$
terhubung), maka $\sigma_A(b) =    \sigma_B(b)$.
 \end{enumerate}
\end{prop}

\begin{proof}[\emph{Petunjuk untuk bukti}] Tinjau fungsi
$f_b\colon (\sigma_A(b))^c \sto B \cap \inv A\colon \lambda \mapsto b - \lambda \vc 1$.   \ns
\end{proof}
